\begin{abstract}
Debiasing methods for scene graph generation are judged by mean recall, and one of the best
known, Total Direct Effect (TDE), raises it on Visual Genome from 14.6 to 24.8. What did
those ten points buy? Mean recall cannot say. A model raises it either by predicting rare
predicates more often within each subject--object pair, or by recognising which individual
images actually show them, and the metric credits both alike. We give an exact way to tell
them apart. Scoring each prediction against the label of a different relation from the same
subject--object group keeps the group's predicate frequencies but breaks the link between
each relation and its own label; this splits any proper-score difference between two frozen models into a
group-level part and a case-level covariance, with no fitted baseline or tuning parameter,
in $O(NK)$ time, with image-clustered confidence intervals. On TDE's released checkpoints
with calibration matched, the proper-score change is almost entirely group-level: under the
quadratic score the case-level part is indistinguishable from zero, and under the log score
it is small and positive, under a tenth of the total. The converse also occurs: a CLIP model
reading image crops, ranked last on accuracy and on the proper score, discriminates better
among relations sharing an object pair than the geometry model ranked above it.
\end{abstract}
